% =============================================================================
%  Notas del profesor — Sesión 12: Torres de Hanoi en Arte ASCII
%  Programación Avanzada — Universidad Panamericana
%  Compilar: pdflatex sesion12_hanoi_ascii_profesor.tex
% =============================================================================
\documentclass[12pt, oneside]{article}
\usepackage[profesor]{progavanzada}

\begin{document}

\portada{Sesión 12}{Torres de Hanoi en Arte ASCII}{2026}

\tableofcontents
\newpage

% =============================================================================
\section{Estructura de la sesión}
% =============================================================================

\begin{center}
\begin{tabular}{clc}
  \toprule
  \textbf{Tiempo} & \textbf{Bloque} & \textbf{Tipo} \\
  \midrule
  10 min & Historia: Lucas 1883, la leyenda del templo & Exposición \\
  10 min & Las tres reglas y $M(n) = 2^n - 1$ & Pizarrón \\
   5 min & Arte ASCII como interfaz — comparación visual & Demo en cuaderno \\
  50 min & Actividad en equipos: construcción del programa & Trabajo en clase \\
   5 min & Cierre: equipos muestran su resultado al grupo & Presentación breve \\
  \midrule
  \textbf{80 min} & & \\
  \bottomrule
\end{tabular}
\end{center}

\begin{notaprofesor}[Contexto de la sesión en el arco de tres clases]
  Esta sesión es el puente entre la \textbf{teoría} (sesión 11, call stack,
  recursión) y la \textbf{práctica recursiva} (sesión 13, algoritmo que
  resuelve Hanoi solo).  El propósito pedagógico es que los alumnos \textit{sientan}
  el problema antes de verlo resuelto por la computadora: si lo mueven a mano,
  la sesión 13 tendrá mucho más impacto cuando la función recursiva lo resuelva
  en milisegundos.

  También es una sesión de \textbf{diseño}: los alumnos eligen cómo se ven las
  torres. Esa diversidad de diseños es insumo para el concurso de la sesión 13.
  Motivarlos a que el código sea legible desde ahora: el equipo ganador tendrá
  que explicarlo frente al grupo.
\end{notaprofesor}

% =============================================================================
\section{Preguntas de arranque}
% =============================================================================

\begin{notaprofesor}[Preguntas para el inicio de clase]
  Lanza estas preguntas antes de abrir el cuaderno:

  \begin{enumerate}
    \item \textbf{``¿Cuántos movimientos harían los sacerdotes del templo de
          Benarés si mueven un disco por segundo, sin parar?''}\\
          \textit{Respuesta esperada:} muchos.  Guiarlos a que calculen
          $2^{64} - 1$ y lo conviertan a años: $\approx 585$ mil millones de años.
          El impacto de ``42 veces la edad del universo'' suele provocar buena reacción.

    \item \textbf{``¿Cuántos movimientos necesito para mover 1 disco? ¿Y para 2?
          ¿Y para 3?''}\\
          \textit{Respuesta:} 1, 3, 7.  Pedir que noten el patrón: siempre
          $2^n - 1$.  No dar la fórmula todavía; dejar que la encuentren ellos.

    \item \textbf{``¿Qué es más difícil: hacer que el programa funcione,
          o hacer que se vea bien en la terminal?''}\\
          \textit{No hay respuesta correcta.} Sirve para motivar la discusión
          sobre el valor del diseño visual en las herramientas de consola.
          Anticipar que algunos alumnos dirán ``lo visual no importa si funciona''
          — usarlo para introducir el argumento de la legibilidad.
  \end{enumerate}
\end{notaprofesor}

% =============================================================================
\section{Notas de contenido ampliado}
% =============================================================================

\subsection{La historia real detrás del juguete}

Édouard Lucas era un matemático serio: contribuyó a la teoría de números,
trabajó con series de Fibonacci y desarrolló pruebas de primalidad.
El juguete de Hanoi fue un proyecto secundario, casi una broma comercial,
pero se convirtió en uno de los objetos matemáticos más citados de la
historia de la computación.

El nombre ``d'Hanoï'' evoca Indochina francesa (Hanói, Vietnam), que en
1883 era colonia francesa.  Lucas eligió el nombre por el exotismo que
evocaba, no por alguna conexión cultural real — del mismo modo que inventó
la leyenda del templo.

\begin{notaprofesor}[Si algún alumno pregunta por la leyenda]
  La historia del templo de Benarés es completamente inventada por Lucas.
  No existe evidencia de ningún templo, ni de ninguna tradición religiosa
  que corresponda a esa descripción.  Lucas la publicó bajo el seudónimo
  ``N. Claus de Siam'' — anagrama de ``Lucas d'Amiens'' — para darle aire
  de misterio oriental.  Un truco de marketing del siglo XIX.
\end{notaprofesor}

\subsection{La conexión binaria: más profunda de lo que parece}

La conexión entre Hanoi y el código Gray no es un mero truco matemático:
es una manifestación de que el \textbf{espacio de estados} del problema
de Hanoi tiene exactamente la misma estructura que el número binario de
$n$ bits.

Cada estado de las torres se puede codificar como un número: para $n=3$
hay $3^3 = 27$ estados (cada uno de los 3 discos puede estar en una de las
3 torres).  Los $2^3 = 8$ estados del camino óptimo corresponden exactamente
a los 8 vértices del cubo de $n$ dimensiones recorrido con el código Gray.

No es necesario explicar esto en clase, pero si algún alumno de matematicas
pregunta, esta es la respuesta completa.

\subsection{¿Por qué es recursivo Hanoi por naturaleza?}

El enunciado del problema ya \textit{es} recursivo: para mover $n$ discos,
necesitas haber movido $n-1$ discos.  No hay forma de describir la solución
sin apelar a la solución del subproblema.  Esto es diferente de problemas
como la suma de una lista, que pueden describirse iterativamente sin rodeo.

La próxima sesión, cuando los alumnos vean el algoritmo recursivo de 6 líneas,
ese ``ajá'' solo funciona si hoy movieron los discos a mano y sintieron la
repetición del patrón.

% =============================================================================
\section{Cómo conducir la actividad en equipos}
% =============================================================================

\begin{notaprofesor}[Formación de equipos]
  Si el grupo no es múltiplo de 3, es preferible un equipo de 4 que dos
  equipos de 2.  Un equipo de 2 tendrá menos capacidad de dividir el
  trabajo (dibujo, movimiento, interfaz) y tiende a producir código menos
  limpio por falta de revisión interna.
\end{notaprofesor}

\subsection{División de roles dentro del equipo}

Sugerir a los equipos que dividan el trabajo así:

\begin{center}
\begin{tabular}{lp{9cm}}
  \toprule
  \textbf{Rol} & \textbf{Responsabilidad} \\
  \midrule
  Diseñador & Decide la forma de los discos, el estilo del poste y la base.
              Trabaja en \texttt{dibujarFila} y \texttt{dibujarTorres}. \\
  Lógica    & Implementa \texttt{moverDisco} con las dos validaciones
              (origen vacío, tamaño). Maneja el estado (\texttt{vector<int>}). \\
  Interfaz  & Escribe el bucle principal, el prompt al usuario, la
              detección de victoria, los mensajes de error. \\
  \bottomrule
\end{tabular}
\end{center}

En la práctica los roles se mezclan, pero tenerlos definidos desde el
principio evita que los tres alumnos escriban código encima del otro.

\subsection{Puntos de control durante los 50 minutos}

\begin{center}
\begin{tabular}{cp{10cm}}
  \toprule
  \textbf{Min} & \textbf{Qué revisar al pasar por los equipos} \\
  \midrule
  15 & ¿Ya definieron el estado como \texttt{vector<int>} con 0 = fondo?
       ¿Ya pueden imprimir una fila de un disco? \\
  30 & ¿\texttt{dibujarTorres} imprime algo coherente?
       ¿El disco del tope se dibuja en la fila correcta? \\
  45 & ¿\texttt{moverDisco} rechaza el movimiento ilegal?
       ¿El bucle pide origen y destino correctamente? \\
  55 & ¿Se detecta la condición de victoria?
       ¿El diseño es legible y consistente? \\
  \bottomrule
\end{tabular}
\end{center}

\begin{notaprofesor}[Equipos bloqueados en el dibujo]
  El problema más frecuente es calcular mal cuántos espacios de relleno
  necesita cada disco.  La fórmula que funciona es:
  \[
    \text{espacios} = \frac{\text{anchoTorre} - (2k - 1)}{2}
  \]
  donde $k$ es el número de disco (1 = más pequeño) y \texttt{anchoTorre}
  es el ancho total de la columna (típicamente 9 u 11).
  Si están usando un carácter diferente a \texttt{=}, la fórmula es la
  misma: el carácter simplemente reemplaza el \texttt{=}.
\end{notaprofesor}

\begin{notaprofesor}[Equipos que quieren implementar el algoritmo recursivo]
  Algunos equipos habrán buscado en internet la solución recursiva de
  Hanoi y querrán implementarla hoy.  \textbf{No permitirlo.}  La
  actividad de hoy es que el \textit{usuario} haga los movimientos.
  Si el programa se resuelve solo, el código del equipo no podrá usarse
  como base para la actividad de la sesión 13 (porque ya la resuelve).
  Decirles que tendrán esa oportunidad la próxima sesión, y que el reto
  de hoy es diferente.
\end{notaprofesor}

% =============================================================================
\section{FAQs frecuentes de los alumnos}
% =============================================================================

\begin{notaprofesor}[FAQ 1 — ¿Por qué \texttt{vector} y no arreglo?]
  \textbf{Pregunta:} ¿Puedo usar un arreglo \texttt{int torre[N]} en lugar
  de \texttt{vector<int>}?\\
  \textbf{Respuesta:} Sí, pero es más incómodo.  Con un arreglo necesitas
  llevar un contador de cuántos discos hay en cada torre.  El
  \texttt{vector} ya tiene \texttt{size()}, \texttt{push\_back()} y
  \texttt{pop\_back()} que modelan exactamente las operaciones de una pila.
  Usar un arreglo no está mal, pero sí innecesariamente complicado.
\end{notaprofesor}

\begin{notaprofesor}[FAQ 2 — ¿Cómo imprimo fila por fila si las torres tienen alturas distintas?]
  \textbf{Pregunta:} Si la torre A tiene 3 discos y la B tiene 1, ¿cómo
  imprimo la fila 2 de B (que está vacía)?\\
  \textbf{Respuesta:} La clave es iterar las filas de arriba a abajo
  (\texttt{fila} desde \texttt{nFilas-1} hasta 0) y, para cada torre,
  verificar si el vector tiene un elemento en esa posición.  Si
  \texttt{fila < torre.size()}, imprime el disco; si no, imprime el poste
  vacío (\texttt{|} centrado con espacios).
\end{notaprofesor}

\begin{notaprofesor}[FAQ 3 — ¿Cuántos discos debo usar?]
  \textbf{Pregunta:} ¿El programa debe soportar cualquier número de discos?\\
  \textbf{Respuesta:} Para la actividad de hoy, con 4 o 5 discos es
  suficiente.  Hardcodear \texttt{int N = 4} no es un problema en esta
  sesión.  En la sesión 13 lo parametrizarán si es necesario.
  Advertir que con muchos discos la terminal se desordena si el ancho
  no está calculado correctamente.
\end{notaprofesor}

\begin{notaprofesor}[FAQ 4 — ¿Cómo muestro el movimiento entre estados?]
  \textbf{Pregunta:} ¿Imprimo las torres antes del movimiento, después, o las dos?\\
  \textbf{Respuesta:} Mínimo una vez después de cada movimiento.
  Algunos equipos querrán limpiar la pantalla (\texttt{system("clear")})
  para que parezca una animación.  Está permitido y puede ser un detalle
  de diseño interesante, pero no es requisito.  Si usan \texttt{system()},
  mencionar que es una llamada al sistema operativo y no funciona igual en
  Windows (\texttt{"cls"}) y en macOS/Linux (\texttt{"clear"}).
\end{notaprofesor}

\begin{notaprofesor}[FAQ 5 — ¿Qué pasa si el usuario escribe una letra en vez de A/B/C?]
  \textbf{Pregunta:} ¿Tenemos que validar que el input sea A, B o C?\\
  \textbf{Respuesta:} Es una buena práctica, pero no es criterio de
  evaluación hoy.  Si lo implementan, sumarles mérito en diseño visual
  (el mensaje de error es parte de la interfaz).  Un \texttt{if (origen
  != 'A' \&\& origen != 'B' \&\& origen != 'C')} es suficiente.
\end{notaprofesor}

% =============================================================================
\section{Evaluación y criterios de calificación}
% =============================================================================

\subsection{Lo que más importa}

Los criterios de evaluación son los de la tabla en las notas del alumno,
pero en la práctica el peso real está en dos cosas:

\begin{enumerate}
  \item \textbf{Que el dibujo sea correcto.} Un programa que no muestra
        las torres de forma coherente no cumple el objetivo fundamental.
        Si el dibujo falla, no importa qué tan bonito sea el carácter
        elegido para los discos.

  \item \textbf{Que la validación rechace discos grandes sobre chicos.}
        Este es el corazón de las reglas. Un programa que acepta cualquier
        movimiento es incorrecto por definición.
\end{enumerate}

\subsection{Criterio de diseño visual — Cómo calibrarlo}

El criterio de diseño (20\%) es el más subjetivo. Usar esta escala:

\begin{center}
\begin{tabular}{cp{9cm}}
  \toprule
  \textbf{Puntos} & \textbf{Descripción} \\
  \midrule
  20 & Los discos se distinguen claramente unos de otros, el poste es
       visible, hay etiquetas A/B/C, la base está bien delimitada,
       el estilo es consistente (mismo carácter en todos los discos). \\
  15 & La torre es legible pero hay inconsistencias menores (e.g., la
       base no está bien alineada o las etiquetas faltan). \\
  10 & Se puede entender qué torre es qué, pero el diseño es mínimo
       o confuso en algunos estados. \\
   5 & El programa imprime algo, pero las torres no son reconocibles
       como torres. \\
  \bottomrule
\end{tabular}
\end{center}

% =============================================================================
\section{Conexión con la sesión 13}
% =============================================================================

\begin{notaprofesor}[No revelar el algoritmo recursivo]
  Al cierre de la sesión, algunos alumnos preguntarán cómo se resuelve
  Hanoi de forma automática.  \textbf{No dar la respuesta.}  Solamente
  decir: ``la próxima sesión verán que el código es sorprendentemente
  corto — por eso importa que el código de hoy esté bien estructurado''.

  El impacto del algoritmo recursivo de sesión 13 depende del contraste
  con el esfuerzo que los alumnos pusieron hoy moviendo los discos a mano.
  Arruinar esa sorpresa reduce el efecto pedagógico.
\end{notaprofesor}

En la sesión 13:

\begin{itemize}
  \item Los equipos presentan su programa (30 segundos cada uno, proyectan
        pantalla o muestran capturas).
  \item El grupo vota por los tres diseños más atractivos con código legible.
  \item El equipo ganador comparte su código en pantalla.
  \item Se agrega la función recursiva \texttt{hanoi(n, origen, destino,
        aux)} al código ganador.
  \item Se ejecuta la solución completa y se verifica que el número de
        movimientos sea $2^n - 1$.
\end{itemize}

Para que la sesión 13 funcione bien, es importante que:

\begin{enumerate}
  \item Los equipos hayan subido su código a Canvas \textbf{antes} de
        que termine la sesión 12 (o como máximo al inicio de la sesión 13).
  \item El código ganador esté disponible en pantalla desde el inicio.
  \item Las tres torres estén representadas como \texttt{vector<int>}
        (o algo equivalente como pilas), de modo que \texttt{moverDisco}
        pueda ser llamada desde la función recursiva sin reescribir.
\end{enumerate}

% =============================================================================
\section{Soluciones a los ejercicios a mano}
% =============================================================================

\begin{notaprofesor}[Ejercicio 1 — 7 movimientos para $n=3$]
  El estado de las torres en cada paso (A, B, C) — lista de discos de
  abajo a arriba, siendo 3 el más grande y 1 el más pequeño:

  \begin{center}
  \begin{tabular}{cccc}
    \toprule
    \textbf{Estado} & \textbf{A} & \textbf{B} & \textbf{C} \\
    \midrule
    0 (inicial) & [3,2,1] & [] & [] \\
    1           & [3,2]   & [] & [1] \\
    2           & [3]     & [2] & [1] \\
    3           & [3]     & [2,1] & [] \\
    4           & []      & [2,1] & [3] \\
    5           & [1]     & [2] & [3] \\
    6           & [1]     & [] & [3,2] \\
    7 (final)   & []      & [] & [3,2,1] \\
    \bottomrule
  \end{tabular}
  \end{center}
\end{notaprofesor}

\begin{notaprofesor}[Ejercicio 2 — $M(4)$ y $M(5)$]
  Verificación: $M(3) = 2^3 - 1 = 7$. \checkmark

  Usando la recurrencia $M(n) = 2\,M(n-1) + 1$:
  \[
    M(4) = 2 \cdot 7 + 1 = 15 \qquad M(5) = 2 \cdot 15 + 1 = 31
  \]
  Verificación directa: $2^4 - 1 = 15$, $2^5 - 1 = 31$. \checkmark
\end{notaprofesor}

\begin{notaprofesor}[Ejercicio 3 — Profundidad del call stack para $n=4$]
  La función recursiva se llama a sí misma con $n-1$, $n-2$, \ldots, 1.
  La profundidad máxima es \textbf{$n = 4$ frames} activos simultáneamente
  en el stack (más el frame de \texttt{main}).  La profundidad del árbol
  de recursión es $n$, no $2^n - 1$ (que es el número de hojas / llamadas totales).
\end{notaprofesor}

\begin{notaprofesor}[Ejercicio 4 — Disco que se mueve cuando $k$ es potencia de 2]
  Cuando $k$ es potencia de 2, su representación binaria tiene exactamente
  un bit en 1: el bit más significativo.  El bit \textit{menos} significativo
  diferente de cero coincide con ese mismo bit.  Por ejemplo:
  $k=4 = 100_2$ → bit 2 → disco 3 (el más grande de la secuencia óptima
  de $n=3$).  En general, cuando $k = 2^j$, se mueve el disco $j+1$.
\end{notaprofesor}

\begin{notaprofesor}[Ejercicio 5 — Desafío: cuatro agujas]
  Con 4 agujas se necesitan \textbf{menos} movimientos.  La intuición:
  la cuarta aguja da un destino temporal adicional, lo que permite
  mover grupos de discos en lugar de uno a la vez.  El resultado exacto
  lo dio Frame en 1941: para $n$ discos y 4 agujas, el mínimo aproximado
  es $\Theta(2\sqrt{n})$ movimientos, mucho menor que $2^n$.
  (No es necesario que los alumnos lleguen a la fórmula — basta la
  intuición de ``menos movimientos porque hay más opciones intermedias''.)
\end{notaprofesor}

% =============================================================================
\section{Resumen para el profesor}
% =============================================================================

\begin{recuerda}
\begin{itemize}
  \item La actividad central son 50 minutos de trabajo en equipos de tres.
  \item El objetivo real no es el código: es que los alumnos \textit{sientan}
        el problema antes de verlo resuelto automáticamente en la sesión 13.
  \item No revelar el algoritmo recursivo — la sorpresa tiene valor pedagógico.
  \item Verificar al minuto 45 que todos los equipos tengan la validación
        de movimiento ilegal funcionando.
  \item Los equipos suben código y capturas a Canvas al final de la sesión.
\end{itemize}
\end{recuerda}

\end{document}
